\subsection{El teorema de Ramanujam--Samuel}
\label{subsection:IV.21.14}
\label{IV.21.14}

\begin{theorem}[21.14.1]
\label{IV.21.14.1}
\textup{(Ramanujam--Samuel)}. ---
Sea $A$ un anillo local noetheriano con ideal maximal $\mathfrak m$, tal que su completado
$\widehat A$ sea íntegro e integralmente cerrado. Sean $B$ un anillo local noetheriano
tal que $\dim(B)\geq\dim(A)$ y $\rho:A\to B$ un homomorfismo local que hace de $B$
una $A$-álgebra formalmente lisa (para las topologías preádicas
$(\mathbf 0, 19.3.1)$), y tal que el cuerpo residual de $B$ sea finito sobre el de $A$.
Entonces todo ciclo $1$-codimensional sobre $\operatorname{Spec}(B)$ que sea principal en el
punto $\mathfrak p=\mathfrak mB$ es un ciclo $1$-codimensional principal.
\end{theorem}

\begin{proof}
Si $k=A/\mathfrak m$ es el cuerpo residual de $A$, $B/\mathfrak mB$ es una
$k$-álgebra formalmente lisa (para su topología preádica) $(\mathbf 0, 19.3.5)$,
\oldpage[IV-4]{324}
por tanto regular y, en particular, íntegra; en otras palabras,
$\mathfrak p=\mathfrak mB$ es efectivamente un ideal primo de $B$, lo que justifica
el enunciado. Basta evidentemente probar que todo ideal primo $\mathfrak q$ de altura $1$\footnote{Corrección de fuente: el francés impreso y el inglés mantenido omiten «de altura $1$». La conclusión concierne a los ciclos de codimensión $1$, y un ideal primo principal no nulo tiene altura $1$.} de
$B$ no contenido en $\mathfrak p$ es \emph{principal}.

Sean $\widehat A$ y $\widehat B$ los completados de $A$ y $B$, respectivamente, de
modo que el ideal maximal de $\widehat A$ es $\mathfrak m\widehat A$; se sabe
$(\mathbf 0, 19.3.6)$ que $\widehat B$ es una $\widehat A$-álgebra formalmente
lisa para las topologías ádicas. Sea $k'$ el cuerpo residual de $\widehat B$,
extensión finita de $k$; existe un homomorfismo local $\widehat A\to C$, donde $C$
es un anillo local noetheriano que es un $\widehat A$-módulo \emph{finito} y
\emph{plano} (por tanto libre), y tal que $C/\mathfrak mC$ es isomorfo a $k'$
$(\mathbf 0_{\mathrm{III}}, 10.3.1)$; de ello se deduce que $C$ es \emph{completo}, y
resulta entonces de (7.5.1), (7.5.3) y (6.5.4, (ii)) que $C$ es \emph{íntegro e
integralmente cerrado}. Además,
\[
  D=\widehat B\widehat\otimes_{\widehat A}C
\]
es un anillo semilocal completo, producto directo de anillos locales completos, uno de los cuales,
$D_0$, tiene por cuerpo residual $k'$ (puesto que $k'\otimes_k k'$ es producto directo
de anillos locales, uno de los cuales es isomorfo a $k'$). Como $D$ es una $C$-álgebra formalmente lisa,
lo mismo ocurre con $D_0$; por consiguiente, $D_0/\mathfrak mD_0$ es una
$k'$-álgebra formalmente lisa, con cuerpo residual $k'$, lo que implica que es
$k'$-isomorfa a un álgebra de series formales
$k'[[T_1,\ldots,T_n]]$ $(\mathbf 0, 19.6.4)$; se concluye, por
$(\mathbf 0, 19.7.1.5)$, que $D_0$ es $C$-isomorfo a
$C[[T_1,\ldots,T_n]]$ y, por tanto, \emph{íntegro e integralmente cerrado} (Bourbaki,
\emph{Alg. comm.}, cap.~V, \S1, n\textsuperscript{o}~4, prop.~14). Como los
morfismos $\operatorname{Spec}(D_0)\to\operatorname{Spec}(\widehat B)$ y
$\operatorname{Spec}(D_0)\to\operatorname{Spec}(B)$ son fielmente planos, se
deduce que $B$ y $\widehat B$ son también \emph{íntegros e integralmente cerrados}
(2.1.13). Esto prueba que el ideal $\mathfrak qD_0$ es divisorial (Bourbaki,
\emph{Alg. comm.}, cap.~VII, \S1, n\textsuperscript{o}~10, prop.~15) y no está contenido
en $\mathfrak mD_0$, pues de lo contrario se tendría
\[
  \mathfrak q=(\mathfrak qD_0)\cap B
  \subset(\mathfrak mD_0)\cap B=\mathfrak mB=\mathfrak p
\]
por planitud fiel $(\mathbf 0_{\mathrm I}, 6.5.1)$. Se concluye que ninguno de los
ideales primos $\mathfrak r_h$ de altura $1$ de $D_0$ que contienen
$\mathfrak qD_0$ puede estar contenido en $\mathfrak mD_0$; si se prueba que estos
ideales son principales, resultará que $\mathfrak qD_0$ es principal, pues los
divisores (en el sentido de Bourbaki) de $\mathfrak qD_0$ y de un producto de potencias de los
$\mathfrak r_h$ son iguales (Bourbaki, \emph{Alg. comm.}, cap.~VII, \S1,
n\textsuperscript{o}~4, prop.~5). Como $\mathfrak qD_0\cap B=\mathfrak q$, se
deduce por planitud fiel que $\mathfrak q$ es principal (2.5.2), utilizando el
hecho de que $B$ es un anillo local.

Quedamos así reducidos a probar el teorema en el caso particular en que $A$ es
\emph{completo} y $B=A[[T_1,\ldots,T_n]]$. Probemos primero que se puede reducir al
caso en que $n=1$. En efecto, procedamos por inducción sobre $n$ y sea
$f\in\mathfrak q$ un elemento que no pertenece a $\mathfrak p=\mathfrak mB$; bastará
probar que existe un $A$-automorfismo $\sigma$ de $B$ tal que
$\sigma(f)$ no pertenece al ideal
\[
  \mathfrak p'=\mathfrak p+BT_1+\cdots+BT_{n-1};
\]
en efecto, si $B'=A[[T_1,\ldots,T_{n-1}]]$, $B'$ es completo y tiene por ideal maximal
$\mathfrak mB'+B'T_1+\cdots+B'T_{n-1}=\mathfrak m'$. Como
$B=B'[[T_n]]$ y $\mathfrak p'=\mathfrak m'B$, resultará de la hipótesis de
inducción (teniendo en cuenta que $B'$ es íntegro e integralmente cerrado) que el ideal
$\sigma(\mathfrak q)$ es principal en $B$, y por tanto también $\mathfrak q$. Ahora bien, si
$\overline f\in k[[T_1,\ldots,T_n]]$ es la imagen de $f$ («serie reducida» de $f$),
se sabe (Bourbaki, \emph{Alg. comm.}, cap.~VII, \S3, n\textsuperscript{o}~7,
lema~3) que existe un $A$-automorfismo $\sigma$ de $B$ tal que
$(\overline{\sigma(f)})(0,\ldots,0,T_n)\neq0$, lo que implica evidentemente
$\sigma(f)\notin\mathfrak p'$.

Supongamos, pues, que $n=1$ y escribamos $T$ en lugar de $T_1$, de modo que $B=A[[T]]$.
Basta probar que, en el anillo de polinomios $A[T]$, el ideal
$\mathfrak q\cap A[T]$ es principal;
\oldpage[IV-4]{325}
en efecto, sea $f_0$ un elemento de $\mathfrak q$ que no pertenece a $\mathfrak mB$;
es una serie formal cuya serie reducida $\overline f_0$ no es nula; por tanto,
en virtud del teorema de preparación (Bourbaki, \emph{Alg. comm.}, cap.~VII, \S3,
n\textsuperscript{o}~8, prop.~5), para todo $f\in\mathfrak q$ existen $g\in B$ y
un polinomio $r\in A[T]$ tales que $f=gf_0+r$, y se tiene entonces
$r\in\mathfrak q\cap A[T]$; por otra parte (\emph{loc. cit.}, prop.~6), existen un
polinomio distinguido no constante $F_0\in A[T]$ y un elemento invertible $u\in B$ tales
que $f_0=uF_0$, por lo que también se tiene $F_0\in\mathfrak q\cap A[T]$, lo que prueba que
$\mathfrak q$ está generado por $\mathfrak q\cap A[T]=\mathfrak q_1$. Puesto que $B$
es plano sobre $A[T]$ $(\mathbf 0_{\mathrm I}, 7.3.3)$, resulta de Bourbaki,
\emph{Alg. comm.}, cap.~VII, \S1, n\textsuperscript{o}~10, prop.~15, que
$\mathfrak q_1$ es un ideal primo de altura $1$ en $A[T]$.

Además, se tiene necesariamente $\mathfrak q_1\cap A=0$; de lo contrario,
$\mathfrak q_1\cap A$ tendría necesariamente altura $\geq1$, y resultaría de
(5.5.3) que $\mathfrak q_1=(\mathfrak q_1\cap A)A[T]$. Pero entonces,
puesto que $\mathfrak q_1\cap A\subset\mathfrak m$, se tendría
$\mathfrak q_1\subset\mathfrak mA[T]$, en contradicción con la hipótesis sobre
$\mathfrak q$. Si $K$ es el cuerpo de fracciones de $A$, $\mathfrak q_1K[T]$ es
entonces un ideal primo distinto de $0$ y de $K[T]$ en $K[T]$, y por tanto de la forma
$hK[T]$, donde
\[
  h(T)=T^m+a_1T^{m-1}+\cdots+a_m,
  \qquad m\geq1,
\]
con $a_i\in K$, siendo $h$ irreducible en $K[T]$. Pero se ha visto más arriba que
$\mathfrak q_1$ contiene un polinomio distinguido no constante $F_0\in A[T]$. Si $t$
es la clase de $T$ en $K[T]/\mathfrak q_1K[T]$, $t$ es, pues, raíz del polinomio
$F_0$ en una extensión de $K$ y, por consiguiente, $h$ divide a $F_0$ en $K[T]$; pero,
puesto que $h$ y $F_0$ son mónicos, ello implica que los coeficientes $a_i$ de $h$
son enteros sobre $A$ (Bourbaki, \emph{Alg. comm.}, cap.~V, \S1,
n\textsuperscript{o}~3, prop.~11), por lo que pertenecen a $A$, ya que $A$ es
integralmente cerrado. En otras palabras, se tiene
\[
  h\in A[T]\cap\mathfrak q_1K[T]=\mathfrak q_1
\]
(Bourbaki, \emph{Alg. comm.}, cap.~II, \S1, n\textsuperscript{o}~5, prop.~11);
como todo polinomio $g\in\mathfrak q_1$ es divisible por $h$ en $K[T]$ y $h$
es mónico, los coeficientes de $g/h$ pertenecen a $A$, por lo que
$\mathfrak q_1=hA[T]$. C.Q.D.
\end{proof}

El enunciado (21.14.1) es \emph{equivalente} al siguiente:

\begin{corollary}[21.14.2]
\label{IV.21.14.2}
Bajo las hipótesis de (21.14.1) sobre $A$, $B$ y $\rho$, para todo ideal primo
$\mathfrak q$ de $B$ no contenido en $\mathfrak p=\mathfrak mB$ y tal que
$\dim(B_\mathfrak q)\geq2$, el anillo $B_\mathfrak q$ es parafactorial. En
particular, si $\dim(B\otimes_A k)>0$ (\emph{i.e.} $\dim B>\dim A$
$((\mathbf 0, 19.7.1)$ y $(\mathbf 0, 6.1.1)))$, el anillo $B$ es parafactorial.
\end{corollary}

\begin{proof}
En efecto, se ha visto en la demostración de (21.14.1) que las condiciones del enunciado
implican que $B$ es íntegro e integralmente cerrado; la equivalencia de los enunciados
(21.14.1) y (21.14.2) resulta entonces de (21.13.15), aplicado a
$X=\operatorname{Spec}(B)$ y $S=\{\mathfrak p\}$.
\end{proof}

\begin{proposition}[21.14.3]
\label{IV.21.14.3}
Sean $S$ un preesquema normal y $f:X\to S$ un morfismo liso.
\begin{enumerate}[label=(\roman*)]
  \item Si $S$ es localmente noetheriano (y por tanto también $X$), entonces, para todo $x\in X$
  tal que $\dim(\mathcal O_{X,x})\geq2$ y que $x$ no sea punto maximal de su
  fibra $f^{-1}(f(x))$, el anillo $\mathcal O_{X,x}$ es parafactorial. Todo ciclo
  $1$-codimensional $Z$ sobre $X$ tal que $\operatorname{Supp}(Z)$ no contenga ninguna
  componente irreducible de una fibra $f^{-1}(s)=X_s$ es localmente principal.
  \item Sea $Y$ un subconjunto cerrado de $X$ que no contenga ninguna componente irreducible
  de una fibra $X_s$, y tal que, para todo punto maximal $\eta$ de $S$,
  $Y_\eta=Y\cap X_\eta$ sea de codimensión $\geq2$ en $X_\eta$. Entonces el par
  $(X,Y)$ es parafactorial.
\end{enumerate}
\end{proposition}

\begin{proof}
Obsérvese que las condiciones de (ii) se cumplen si $Y$ es un subconjunto cerrado de $X$
tal que, para todo $s\in S$, $Y_s=Y\cap X_s$ sea de codimensión $\geq2$ en
$X_s$.

Para probar (21.14.3), observemos primero que, bajo las hipótesis de (i), si se pone
$Y=\overline{\{x\}}$, existe un entorno abierto $V$ de $x$ en $X$ tal que
$Y\cap V$ y $V$ satisfacen las condiciones de (ii): en efecto, resulta de la hipótesis
y de (9.5.3) que se puede tomar $V$ de modo que $V\cap Y$ no contenga ninguna
componente irreducible de un $X_s$; por otra parte, para todo punto maximal $\eta$ de
$S$ tal que $Y\cap X_\eta\neq\varnothing$, los puntos de $Y\cap X_\eta$ son
especializaciones de $x$, por lo que en un punto tal $z$ se tiene
$\dim(\mathcal O_{X,z})\geq2$ y, como
$\mathcal O_{X,z}=\mathcal O_{X_\eta,z}$ ($S$ es reducido en el punto $\eta$ en virtud
de (2.1.13)), $Y\cap X_\eta$ es efectivamente de codimensión $\geq2$ en $X_\eta$. Como el
par $(V,Y\cap V)$ es entonces parafactorial en virtud de (ii), se concluye de (21.13.10)
que el anillo $\mathcal O_{X,x}$ es parafactorial, es decir, la primera afirmación
de (i). Para probar la segunda, se puede reducir al caso en que
$Z=\overline{\{z\}}$, donde $z\in X^{(1)}$; como $\mathcal O_{X,z}$ es un anillo local
integralmente cerrado y de dimensión $1$, es un anillo de valoración discreta y $Z$
es, por tanto, principal en el punto $z$; para todo otro punto $x$ de
$\operatorname{Supp}(Z)$, se tiene entonces $\dim(\mathcal O_{X,x})\geq2$ y, por hipótesis, $x$ no es
punto maximal de su fibra, por lo que $\mathcal O_{X,x}$ es parafactorial.
Apliquemos entonces (21.13.15), tomando como $S$ el conjunto formado por $z$ y los puntos
maximales de las fibras de $f$; es claro que $Z$ es principal en los puntos de $S$
puesto que $z$ es el único punto de $S$ que pertenece a $\operatorname{Supp}(Z)$; como,
por otra parte, la condición $b)$ de (21.13.15) se cumple evidentemente, se concluye que
$Z$ es localmente principal.

Quedamos, pues, reducidos a probar (ii). Pongamos $U=X-Y$ y sea $j:U\to X$ la inyección
canónica.
\oldpage[IV-4]{326}
Como las hipótesis y la conclusión son locales sobre $S$ y sobre $X$ en virtud de
(21.13.6, (i)), se puede reducir al caso en que $S$ y $X$ son afines, siendo entonces $f$ de
presentación finita. Puesto que las fibras de $f$ son regulares, la hipótesis sobre $Y$
implica que, para todo punto $x\in Y$, se tiene
\[
  \operatorname{prof}(\mathcal O_{X_{f(x)},x})
  =\dim(\mathcal O_{X_{f(x)},x})\geq1;
\]
resulta de (19.9.8) que el homomorfismo canónico
$\mathcal O_X\to j_*(\mathcal O_U)$ es \emph{inyectivo}. Además, sustituyendo $Y$ por
un subconjunto cerrado mayor según el método de la demostración de (19.9.8), y
utilizando (21.13.6, (ii)), se puede suponer que $Y$ está definido por un ideal \emph{de
tipo finito} del anillo de $X$; por tanto, el abierto $U=X-Y$ es \emph{cuasicompacto y
cuasiseparado} y el conjunto cerrado $Y$ es \emph{constructible}.

Para probar (ii), basta, dados un $\mathcal O_U$-módulo invertible $\mathscr L$
y un punto $x\in Y$, establecer la existencia de un entorno abierto $V$ de $x$
en $X$ tal que: $1^{\mathrm o}$ el homomorfismo canónico
$\mathcal O_V\to(j_V)_*(\mathcal O_{U\cap V})$ sea sobreyectivo;
$2^{\mathrm o}$ exista un $\mathcal O_V$-módulo invertible $\mathscr L'_V$ tal que
$\mathscr L'_V|(U\cap V)=\mathscr L|(U\cap V)$ ((21.13.5) y (21.13.3)).

Pongamos $s=f(x)$; veremos que se puede reducir al caso en que
$S=S_0=\operatorname{Spec}(\mathcal O_{S,s})$. Sea, en efecto, $(S_\nu)$ un sistema
fundamental de entornos abiertos afines de $s$ en $S$, y pongamos
$X_\nu=f^{-1}(S_\nu)$, $Y_\nu=Y\cap X_\nu$, $U_\nu=X_\nu-Y_\nu$,
siendo $j_\nu:U_\nu\to X_\nu$ la inyección canónica; $X_0=X\times_S S_0$ es entonces
límite proyectivo de los $X_\nu$ (8.1.2, $a$)); supongamos verdadera la proposición para el
morfismo $f_0:X_0\to S_0$ y para $Y_0=Y\cap X_0$, y pongamos $U_0=X_0-Y_0$,
siendo $j_0:U_0\to X_0$ la inyección canónica. Como la proyección
$p_\nu:X_0\to X_\nu$ es un morfismo afín, se tiene
\[
  (j_0)_*(\mathcal O_{U_0})
  =p_\nu^*((j_\nu)_*(\mathcal O_{U_\nu}))
  \qquad(\mathbf{II}, 1.5.2)
\]
y el homomorfismo canónico
$\rho_0:\mathcal O_{X_0}\to(j_0)_*(\mathcal O_{U_0})$ no es otro que
$p_\nu^*(\rho_\nu)$, donde
$\rho_\nu:\mathcal O_{X_\nu}\to(j_\nu)_*(\mathcal O_{U_\nu})$ es el homomorfismo
canónico; como, por hipótesis, $\rho_0$ es sobreyectivo, lo mismo ocurre con $\rho_\nu$
para $\nu$ suficientemente grande (8.5.7). Sea, por otra parte, $\mathscr L_0$ el
$\mathcal O_{U_0}$-módulo invertible obtenido por restricción de $\mathscr L$, y observemos que los
$X_\nu$ son afines, y por tanto cuasicompactos y cuasiseparados; puesto que, por hipótesis,
existe un $\mathcal O_{X_0}$-módulo invertible $\mathscr L'_0$ tal que
$\mathscr L'_0|U_0=\mathscr L_0$, existen un $\nu$ suficientemente grande y un
$\mathcal O_{X_\nu}$-módulo cuasicoherente $\mathscr L'_\nu$ de presentación finita
tales que $\mathscr L'_0$ sea isomorfo a $p_\nu^*(\mathscr L'_\nu)$ (8.5.2, (ii));
además, por (8.5.5), se puede suponer que $\mathscr L'_\nu$ es invertible. Finalmente, como
los $U_\nu$ son cuasicompactos y cuasiseparados y
$\mathscr L'_0|U_0=\mathscr L_0$, existe un $\nu$ suficientemente grande tal que
$\mathscr L'_\nu|U_\nu$ y $\mathscr L|U_\nu$ sean isomorfos (8.5.2.5).

\oldpage[IV-4]{327}
Quedamos así reducidos al caso en que $S=\operatorname{Spec}(A)$,
$X=\operatorname{Spec}(B)$, donde $A$ es un anillo \emph{local}; puesto que $X$ es
normal y $f$ liso, $S$ es normal (17.5.7), y por tanto $A$ es \emph{íntegro e
integralmente cerrado}. Consideremos entonces $A$ como límite inductivo de sus
\emph{sub-$\mathbb Z$-álgebras de tipo finito}; como la clausura integral de una
sub-$\mathbb Z$-álgebra de este tipo es un subanillo de $A$ y es también una
$\mathbb Z$-álgebra de tipo finito (7.8.3, (ii), (iii) y (vi)), $A$ es límite inductivo
filtrante de las sub-$\mathbb Z$-álgebras de tipo finito $A_\lambda$ de $A$ que son
anillos \emph{integralmente cerrados}. En virtud de (1.8.4.2), existen un índice
$\lambda$ y una $A_\lambda$-álgebra de tipo finito $B_\lambda$ tales que
\[
  B=B_\lambda\otimes_{A_\lambda}A
\]
salvo $A$-isomorfismo. Pongamos $S_\lambda=\operatorname{Spec}(A_\lambda)$,
$X_\lambda=\operatorname{Spec}(B_\lambda)$; si $p_\lambda:X\to X_\lambda$ es la
proyección canónica, se puede además suponer (puesto que $Y$ es
\emph{constructible}) que $Y=p_\lambda^{-1}(Y_\lambda)$, donde $Y_\lambda$ es un
subconjunto cerrado de $X_\lambda$ (8.3.11). Sea $f_\lambda$ el morfismo
$X_\lambda\to S_\lambda$; en virtud de la transitividad de las fibras y de (4.2.6),
$Y_\lambda$ no contiene ninguna componente irreducible de las fibras
$f_\lambda^{-1}(s_\lambda)$ para ningún $s_\lambda\in S_\lambda$. Puesto que $f$ es
liso, se puede también suponer que $f_\lambda$ es liso (17.7.8); finalmente, la imagen del
punto genérico $\eta$ de $S$ en $S_\lambda$ es el punto genérico $\eta_\lambda$
de $S_\lambda$, y se tiene
\[
  \operatorname{codim}((Y_\lambda)_{\eta_\lambda},
  (X_\lambda)_{\eta_\lambda})
  =\operatorname{codim}(Y_\eta,X_\eta)\geq2
\]
en virtud de (6.1.4). Se ve, pues, que $S_\lambda$, $X_\lambda$ y $Y_\lambda$
satisfacen todas las hipótesis de (ii), y, poniendo
$X_\mu=X_\lambda\times_{S_\lambda}S_\mu$ para $\lambda\leq\mu$, lo mismo ocurre con
$S_\mu$, $X_\mu$ y $Y_\mu$. Probemos que, si se demuestra que el par
$(X_\mu,Y_\mu)$ es parafactorial para todo $\mu\gg\lambda$, lo mismo ocurre con
$(X,Y)$. En efecto, sean $U=X-Y$, $U_\mu=X_\mu-Y_\mu$, y sean $j:U\to X$ y
$j_\mu:U_\mu\to X_\mu$ las inyecciones canónicas; como la proyección
$p_\mu:X\to X_\mu$ es un morfismo afín, se tiene
\[
  j_*(\mathcal O_U)=p_\mu^*((j_\mu)_*(\mathcal O_{U_\mu}))
  \qquad(\mathbf{II}, 1.5.2),
\]
y por consiguiente el homomorfismo canónico
$\rho:\mathcal O_X\to j_*(\mathcal O_U)$ no es otro que $p_\mu^*(\rho_\mu)$, donde
$\rho_\mu:\mathcal O_{X_\mu}\to(j_\mu)_*(\mathcal O_{U_\mu})$ es el homomorfismo
canónico; como este último es biyectivo por hipótesis, lo mismo ocurre con $\rho$.
Por otra parte, para todo $\mathcal O_U$-módulo invertible $\mathscr L$, existen un
$\mu\gg\lambda$ y un $\mathcal O_{U_\mu}$-módulo invertible $\mathscr L_\mu$ tales
que $\mathscr L=p_\mu'^*(\mathscr L_\mu)$ (denotando por
$p_\mu':U\to U_\mu$ la restricción de $p_\mu$) (8.5.2 y 8.5.5), pues $U_\lambda$ es
cuasicompacto ya que $X_\lambda$ es noetheriano. Por hipótesis, existe un
$\mathcal O_{X_\mu}$-módulo invertible $\mathscr L'_\mu$ tal que
$\mathscr L'_\mu|U_\mu$ sea isomorfo a $\mathscr L_\mu$;
$\mathscr L'=p_\mu^*(\mathscr L'_\mu)$ es entonces un $\mathcal O_X$-módulo
invertible tal que $\mathscr L'|U$ sea isomorfo a $\mathscr L$, lo que prueba nuestra
afirmación.

Quedamos así reducidos a probar (ii) cuando el anillo $A$ es una
$\mathbb Z$-álgebra de tipo finito integralmente cerrada; como los anillos locales
$\mathcal O_{S,s}$ de $S$ son entonces anillos excelentes integralmente cerrados, sus
\emph{completados} son también integralmente cerrados (7.8.3, (ii), (iii) y (vii)). Para
probar que el par $(X,Y)$ es parafactorial, basta mostrar que, para todo
$x\in Y$, el anillo $\mathcal O_{X,x}$ es parafactorial (21.13.10). Sea $y$ un punto
cerrado de $Y_{f(x)}$ que sea especialización de $x$ (5.1.11); si se pone
$s=f(x)=f(y)$, $\mathcal O_{S,s}$ tiene un
\oldpage[IV-4]{328}
completado integralmente cerrado, y $\mathcal O_{X,y}$ es una
$\mathcal O_{S,s}$-álgebra formalmente lisa para las topologías preádicas
(17.5.3); puesto que $Y$ no contiene ninguna componente irreducible de $X_s$, se tiene
$\dim(\mathcal O_{X,y})>\dim(\mathcal O_{S,s})$ (17.5.8). Si $s\neq\eta$, se tiene
$\dim(\mathcal O_{S,s})\geq1$; si, por el contrario, $s=\eta$, se tiene por hipótesis
$\dim(\mathcal O_{X,x})\geq2$; por tanto, en todos los casos,
$\dim(\mathcal O_{X,y})\geq2$. Además, el ideal primo de $\mathcal O_{X,y}$
correspondiente al punto $x$ no está contenido en
$\mathfrak m_s\mathcal O_{X,y}$, puesto que $Y$ no contiene ninguna componente
irreducible de $X_s$. Finalmente, puesto que $y$ es cerrado en $X_s$, $k(y)$ es una extensión
finita de $k(s)$ $(\mathbf I, 6.4.2)$. En todos los casos, se puede aplicar a
$A=\mathcal O_{S,s}$, $B=\mathcal O_{X,y}$ y
$B_\mathfrak q=\mathcal O_{X,x}$ el resultado de (21.14.2), lo que termina la
demostración.
\end{proof}

\begin{env}[21.14.4]
\label{IV.21.14.4}
\emph{Observaciones}. ---

(i) Es posible que los enunciados (21.14.1) y (21.14.2) sigan siendo válidos \emph{sin
hipótesis} sobre el cuerpo residual de $B$. Formulado con esta generalidad, el resultado
sería equivalente, en virtud de (21.13.15) y (21.13.12.1), al siguiente: sean $A$ un
anillo local noetheriano completo, íntegro e integralmente cerrado, y $B$ un anillo local
noetheriano que sea una $A$-álgebra formalmente lisa, tal que
$\dim(B)>\dim(A)$; entonces $B$ es parafactorial.

(ii) Veremos en el cap.~VI, empleando una técnica de «descenso finito»
aplicada al morfismo $Y'\to Y=\operatorname{Spec}(A)$, donde $Y'$ es la normalización de
$Y$, que la conclusión de (21.14.1) (o de (21.14.2)) sigue siendo válida si se sustituye
la hipótesis de que $\widehat A$ es integralmente cerrado por la hipótesis de que
$\widehat A$ es \emph{reducido}, siempre que $\dim B\geq\dim A+2$. Si se sustituye esta
última condición por $\dim B\geq\dim A+3$, se puede incluso suprimir la hipótesis de que
$\widehat A$ es reducido. Del mismo modo, la conclusión de (21.14.3) sigue siendo válida si se
sustituye la hipótesis de que $X$ es normal por la hipótesis de que $X$ es reducido, siempre
que $\dim(\mathcal O_{X,x})\geq\dim(\mathcal O_{S,f(x)})+2$.

(iii) En el cap.~III, 3\textsuperscript{a} parte, se prueba que, si $f:X\to S$ es un
morfismo liso y $Y$ es un subconjunto cerrado localmente constructible de $X$ tal que,
para todo $s\in S$, se tenga (con las notaciones de (21.14.3))
$\operatorname{codim}(Y_s,X_s)\geq3$, entonces el par $(X,Y)$ es parafactorial
([41], XII, 4.8). Esta conclusión deja de ser válida cuando solo se supone que
$\operatorname{codim}(Y_s,X_s)\geq2$ para todo $s\in S$ y ya no se supone que
$X$ sea reducido. Por ejemplo, sean $k$ un cuerpo y $A=k[T]/(T^2)$, el álgebra de los números
duales sobre $k$, $S=\operatorname{Spec}(A)$,
$X=\operatorname{Spec}(A[T_1,T_2])$ ($T_1$, $T_2$ indeterminadas), de modo que
$X=\mathbb V_S^2$, siendo $Y$ la «sección nula» de este fibrado; si $s$ es el único
punto cerrado de $S$, se tiene $X_s=\mathbb V_k^2$ y $Y_s$ es la «sección nula» de
$X_s$. Para ver que el par $(X,Y)$ no es parafactorial, basta probar
que el anillo $B=A[T_1,T_2]_\mathfrak m$, donde $\mathfrak m$ es el ideal
$(T_1)+(T_2)$ que define $Y$, no es parafactorial (21.13.10). Sea
$B_0=B_{\mathrm{red}}$, y denotemos por $U$ y $U_0$ los complementos del punto cerrado
en $\operatorname{Spec}(B)$ y $\operatorname{Spec}(B_0)$; razonando como
en (21.13.9), se tiene la sucesión exacta, extensión de (21.13.9.1),
\[
  \Gamma(U,\mathcal O_U)^*\longrightarrow\Gamma(U_0,\mathcal O_{U_0})^*
  \longrightarrow\mathrm H^1(U_0,\mathcal O_{U_0})
  \longrightarrow\operatorname{Pic}(U)\longrightarrow\operatorname{Pic}(U_0).
\]
Ahora bien, se tiene $\Gamma(U,\mathcal O_U)=B$, $\Gamma(U_0,\mathcal O_{U_0})=B_0$ puesto que
$\operatorname{prof}(B)=\operatorname{prof}(B_0)=2$ (5.10.5). Por otra parte,
$\operatorname{Pic}(U_0)=0$ puesto que $B_0$ es factorial, y
$\mathrm H^1(U_0,\mathcal O_{U_0})\neq0$ [41, 3, Ejemplo III-1]; como
el homomorfismo $B^*\to B_0^*$ es sobreyectivo, se concluye que
$\operatorname{Pic}(U)\neq0$, y por tanto que $B$ no es parafactorial.
\oldpage[IV-4]{329}
(iv) El resultado (21.14.3) fue demostrado por primera vez por C.~Seshadri [44] en el caso
particular en que $S$ es un esquema algebraico normal sobre un cuerpo algebraicamente cerrado
$k$ y $X=S\times_k T$, donde $T$ es un $k$-preesquema algebraico liso sobre $k$. La
demostración de Seshadri [44, p.~188--189] es de naturaleza global y utiliza la
teoría de los esquemas de Picard. Proporciona además (\emph{loc. cit.}) resultados
como el siguiente (para el que actualmente no se dispone de una demostración
por métodos locales). Sean $S$ y $T$ dos preesquemas localmente de tipo finito sobre un
cuerpo $k$, $X=S\times_k T$, y $Z$ un ciclo $1$-codimensional sobre $X$ (considerado como
$S$-preesquema); supongamos que se cumplen las condiciones siguientes:
\begin{enumerate}[label=\arabic*\textsuperscript{o}]
  \item $S$ y $T$ son geométricamente normales sobre $k$ (6.7.6);
  \item Para todo punto maximal $\eta$ de $S$, el ciclo $1$-codimensional $Z_\eta$
  sobre la fibra $X_\eta$, que tiene la misma multiplicidad que $Z$ en todo punto de
  $X^{(1)}\cap X_\eta$, es localmente principal (en otras palabras, es la imagen de un
  divisor de $X_\eta$, puesto que $X_\eta$ es normal);
  \item Para todo $s\in S$, $Z$ es principal en los puntos maximales de la fibra $X_s$.
\end{enumerate}
Entonces $Z$ es localmente principal. En otras palabras, como $X$ es normal (6.14.1),
para todo $x\in X$ que no sea maximal en su fibra $X_s$ y que no pertenezca a
ninguna de las «fibras genéricas» $X_\eta$ (lo que implica
$\dim(\mathcal O_{X,x})\geq2$ en virtud de (6.1.1)), el anillo local
$\mathcal O_{X,x}$ es parafactorial, en virtud de (21.13.15)\footnote{Corrección de referencia: el francés impreso y el inglés mantenido dan (21.12.15), que trata de un criterio de isomorfía para un morfismo propio; el criterio de parafactorialidad pertinente es (21.13.15).}.
\footnote{En realidad, en el artículo citado, Seshadri supone que $k$ es algebraicamente
cerrado, que $T$ es separado y «semicompleto» (i.e.\ tal que
$\Gamma(T,\mathcal O_T)$ sea $k$-isomorfo a $k$), y sustituye la hipótesis
$3^{\mathrm o}$ por la hipótesis más fuerte de que $\operatorname{Supp}(Z)$ no contiene
ninguna de las fibras $X_s$ para $s\in S$. Pero, como el enunciado es local sobre $S$, se
concluye inmediatamente que basta formular la hipótesis $3^{\mathrm o}$, y esto prueba que
la conclusión (interpretada como antes en términos de la propiedad de parafactorialidad
de los anillos $\mathcal O_{X,x}$) es local sobre $S$ y sobre $T$, lo que permite
eliminar por completo la hipótesis de que $T$ es «semicompleto» y de que $k$ es
algebraicamente cerrado, puesto que (después de pasar, si es necesario, a la clausura algebraica de $k$)
se puede suponer inicialmente que $T$ es afín, lo que permite sumergirlo como abierto de un
esquema proyectivo y normal sobre $k$, al que se aplica el resultado de Seshadri. Obsérvese
también que, gracias a esta reducción, basta hacer la demostración de Seshadri en
el caso en que $T$ es \emph{proyectivo} (y no solo «semicompleto»), caso en el que la
teoría del esquema de Picard utilizada por Seshadri está contenida en la que se
desarrollará en el cap.~VI de nuestro Tratado.}
\end{env}
